\section{INTRODUÇÃO}

Louças sanitárias combinam massa elevada, geometria irregular e comportamento frágil. O produto percorre várias interfaces até a instalação, e o ponto em que o dano se torna visível pode não coincidir com o momento em que a falha surgiu. Essa diferença dificulta a atribuição de causa, distribui perdas entre empresas e trabalhadores e pode levar a respostas locais que tratam apenas o efeito.

O material do projeto organiza a cadeia em fabricação e distribuição, transporte externo, recebimento e abertura, movimentação interna e instalação. A equipe observou o contexto de recebimento, estoque e movimentação em uma obra residencial em acabamento e ouviu atores envolvidos no fluxo. Os relatos convergiram para uma avaria anterior à instalação, com o transporte do distribuidor até a obra como recorte prioritário \citep{equipe2026}.

O objetivo deste relatório é documentar a definição do problema, explicitar as evidências que sustentam o recorte e registrar progressivamente as etapas padronizadas de desenvolvimento adotadas na disciplina. O trabalho evita antecipar uma concepção de produto e preserva a sequência entre necessidade, requisito do cliente, requisito mensurável, função, princípio de solução e concepção \citep{rozenfeld2006,back2008}.

Esta versão representa uma entrega intermediária. O conteúdo está organizado nas mesmas etapas do relatório-base: pré-desenvolvimento, projeto informacional, projeto conceitual e consideração final. As afirmações provenientes da observação e dos relatos são qualitativas e serão progressivamente substituídas ou complementadas por evidências rastreáveis, indicadores e critérios de decisão.
